\section{Conclusion}

\label{sec:conclusion}

This volume has traced the structure
of divergence, renormalization,
and the limits of modern physical theory.

It has introduced the \DH{}
as a boundary between the \generated{}
and the \ungenerated{}.

\medskip
\begin{quote}
It does not explain what lies beyond,\\
but clarifies where explanation ceases.
\end{quote}
\medskip

The \ungenerated{} domain
cannot be described.

It is not an object of theory,
but appears only as a boundary.

\medskip
\begin{quote}
The boundary is not an object,\\
but a condition.
\end{quote}
\medskip

Divergence appears
as theory approaches this boundary.

Renormalization converts that contact
into a form that can be handled within theory.

These are not separate phenomena.

\medskip
\begin{quote}
They are structural responses\\
to the same boundary.
\end{quote}
\medskip

Several questions remain open.

The precise mathematical structure of the \DH{}---
its thickness, fluctuation spectrum,
and relation to known geometric structures---
awaits further development.

The connection between this framework
and existing approaches to quantum gravity,
holography, and information-theoretic methods
remains to be made explicit.

The \DH{} is not a place that can be reached.

It marks the limit
of what can be described.

\medskip
\begin{quote}
Nothing further can be said.
\end{quote}
